Las flores, frutas, vegetales y hongos representan un pilar fundamental en la cadena alimentaria \citep{Cassani_2022}. La industria agrícola, responsable de su producción continua, enfrenta desafíos crecientes debido a la incertidumbre respecto a las condiciones de cultivo, las cuales se han visto afectadas, en gran parte, por el cambio climático \citep{Parajuli_2019}.\\

Las alteraciones que se derivan del cambio climático repercuten en diversos aspectos, entre ellos, el ciclo de producción agrícola, debido a variaciones inusuales e inesperadas en temperatura, precipitación y patrones de migración de polinizadores, entre otros factores. Los cuales, en conjunto, han modificado las zonas, temporadas y rendimiento de los cultivos \citep{Bayable_2021, Cuevas_2021, Parajuli_2019}.\\

La estabilidad de los procesos productivos depende de cierta consistencia en las condiciones ambientales. Sin embargo, las variaciones a las que nos enfrentamos actualmente, derivadas principalmente por el cambio climático, repercuten en dicha consistencia, incrementando la inestabilidad en la seguridad alimentaria y el desabasto a la par de reducir la producción y exportación, impactando directamente la economía de las naciones \citep{Farooq_2022}.\\

Cada cultivo prospera dentro de umbrales agroambientales específicos, que en muchos casos, han sido ya ampliamente descritos con anterioridad en la literatura científica; e.g., el aguacate \citep{Moraes_2022, Rocha-Arroyo_2011} o el agave \citep{García-Mendoza_2002}. El maíz (\textit{Zea mays}), un cultivo de origen mexicano \citep{OrozcoRamrez_2016}, representa un caso particular al ser considerado como uno de los cultivos más importantes a nivel mundial para el abastecimiento alimenticio directo e indirecto y sus usos industriales \citep{Revilla_2022}.\\

El maíz ha sido empleado para consumo tanto de dieta primaria como de dieta secundaria; i.e., incorporado en conjunto con otros alimentos o para forrajeo \citep{Revilla_2022}. Asimismo, este cultivo se emplea en diversas actividades del sector productivo y de manufactura \citep{Ranum_2014}; e.g., cervezas, jarabes, harinas, medicamentos y biocombustibles \citep{Rausch_2006,SOKANADEAGA_2024, Zhang_2022}.\\


A nivel mundial, el consumo de las distintas variedades de este cultivo aporta, en promedio, el $5.3\%$, el $4.6\%$ y $1.6\%$ de la ingesta calórica, proteica y grasa, respectivamente, per capita diaria \citep{Erenstein_2022}. Las variedades del maíz pueden clasificarse por su tamaño, dulzor, composición del endospermo, color o tasa de maduración \citep{Ranum_2014}; no obstante, comparten su respuesta a las condiciones de cultivo \citep{Baum_2019, Sanchez_2014, Parent_2012}.\\


Los calendarios agrícolas proporcionan información crítica sobre cuándo se llevan a cabo las etapas de desarrollo del cultivo que conforman su fenología \citep{Franch_2022}. De esta manera, se organiza el manejo agrícola tanto de la secuencia fenológica completa como de etapas particulares; e.g., la siembra o la cosecha.\\

A través del alza en la frecuencia de eventos climáticos extremos \citep{MIRON_2023}, el incremento de la temperatura promedio del planeta \citep{Shivanna_2022} y los desfases en la duración de las estaciones del año \citep{Wang_2014}, el cambio climático ha alterado la distribución espacial de las zonas de cultivo \citep{SU_2023} y los calendarios agrícolas \citep{Franch_2022}.\\


Lo anterior impone a la agricultura un reto para garantizar la calidad y la cantidad de productos para satisfacer las necesidades alimentarias en el futuro \citep{MIRON_2023}. Por otro lado, se ha observado que, para el maíz, las fechas de siembra juegan un rol determinante en la cosecha resultante \citep{Baum_2019, BernalMorales_2025, Long_2017, Mason_2018, Parker_2016, Sacks_2011}.\\


De acuerdo con \cite{Sanchez_2013}, el periodo de la siembra a la emergencia de la plántula es una de las temporadas más susceptibles de este cultivo. Por tanto, este trabajo se centra en proponer una recalendarización del periodo de siembra que garantice una baja exposición al estrés ambiental durante el ciclo de desarrollo comprendido hasta la emergencia.\\

Anualmente se destinan aproximadamente 200 millones de hectáreas (M ha) para la plantación de maíz, de las cuales, aproximadamente el 9\% se ubica en Europa \citep{Erenstein_2022}. No obstante, a consecuencia del cambio climático, las regiones con mayor altitud y aquellas al norte de la región noroeste de Europa, en particular Alemania, tendrán un incremento en las zonas con potencial de producción agrícola del maíz \citep{Miedaner_2021}.\\

Esta transición implica retos en tanto a la renovación de la planeación, la gestión, el desarrollo y la correcta implementación de las técnicas agrícolas relativas al manejo y cuidado del cultivo \citep{Parent_2018, DONMEZ_2023}. La identificación de periodos de siembra, bajo la transición que supone el cambio climático, continúa siendo limitada, no por desinterés sino por la tendencia a considerarlas estáticas \citep{Parker_2016}, la complejidad inherente a la incertidumbre climática \citep{Waha_2011} y por falta de especificidad en la información disponible correspondiente a las fechas de plantación  \citep{Portmann_2010, Sacks_2010}.\\

 En este sentido, este estudio pretende robustecer la planificación de las fechas de siembra del maíz mediante la implementación de herramientas analíticas y computacionales para la toma de esta decisión, ante la necesidad de adaptar la planificación agrícola a los cambios esperados en las condiciones ambientales futuras. El estudio se focaliza en el sureste de Alemania, una región que alberga aproximadamente el 54.42\% de las 9.5 M ha de superficie empleada anualmente para la producción de los principales cultivos \citep{Duden_2024}.\\

La expansión proyectada de su área potencial de cultivo en regiones como Alemania, la estabilidad nutricional bajo escenarios de elevadas concentraciones atmosféricas de CO$_2$ \citep{Myers_2014} y su papel en la seguridad alimentaria global \citep{TANUMIHARDJO_2020}, en conjunto, consolidan al maíz como un cultivo estratégico no solo por su relevancia productiva y económica ante el cambio climático.\\


En consecuencia, este trabajo busca cubrir la necesidad de actualizar la planificación agronómica del periodo comprendido de la siembra a la emergencia. Determinando, mediante el análisis de datos agroclimáticos históricos, aquellas temporadas favorables para el desarrollo, mitigando el efecto del cambio climático.
